\documentclass[a4paper,11pt]{article}

\usepackage{revy}
\usepackage[utf8]{inputenc}
\usepackage[T1]{fontenc}
\usepackage[danish]{babel}

\revyname{Matematikrevyen}
\revyyear{2026}
\version{0.1}
\eta{$4$ minutter}
\status{Færdig}

\title{AI i aflevering}
\author{Niklas '21 \& Thais '22}

\begin{document}
\maketitle

\begin{roles}
\role{S}[Skuespiller] "Studerende"
\role{T}[Skuespiller] "TA"
\role{U}[Skuespiller] "Underviser"
\end{roles}

\begin{sketch}
\scene{S sidder i et afhøringslokale og T kommer ind med papirer i hånden.}
\says{T} Bare så det er helt klart S. Jeg er fuldstændig ligeglad med at du er tidligere Georg Mohr vinder. Hvad jeg ikke er ligeglad med, er at jeg 100 procent ved, at du har brugt AI til at lave din aflevering.
\scene{T slår papirerne, som er afleveringen, på bordet.}
\says{T} Bare indrøm det.
\says{S} Måske... men du har ingen beviser.
\says{T} Det er jeg faktisk glad for at du siger, for der er jeg uenig.
\scene{T bladrer i afleveringen og begynder at læse op}
\says{T} "Da vi tydeligvis har en normal familie af holomorfe funktioner ved vi per Montels theorem, at familien er lokal uniform bounded, derfor er den konstante funktion f(x)=5, begrænset". Det her er matintro, du kender slet ikke Montels theorem, det er fra et kandidatkursus.
\says{S} Det betyder ingenting. Resultatet... kom til mig i en drøm. Desuden, så ved jeg ikke engang hvad det her "AI", er for noget.
\says{T} Du ved ikke hvad AI er? Det har jeg virkelig svært ved at tro på.
\scene{T læser igen op}
\says{T} "Note til mig selv, spørg chatten om det her senere". Bare indrøm det.
\says{S} Jeg indrømmer at chatten er kælenavnet på min studiekammerat.
\scene{T kigger helt vantro og målløs}
\says{S} Det er bare definitioner, lemmaer og theoremer, der bare er blevet skrevet i en dope ass rækkefølge. Men du har stadigvæk intet bevis.
\scene{T bladrer videre og vender papirerne mod S og peger på dem.}
\says{T} "Chat, hvordan løser jeg spørgsmål 3.a, 3.b og 3.c? Gør det lidt fancy så TA bliver imponeret"... Du har inkluderet en prompt i din aflevering!
\says{S} Jeg aner ikke hvad det er, det er første gang jeg ser det.
\scene{T bliver mere frustreret}
\says{S} Jeg er ked af at sige det T, men jeg brugte Montel fordi det er det jeg føler var den åbenlyse løsning, da alle tal bare specielt kan ses som værende konstante funktioner, som derfor er holomorfe og familien er normal.
\scene{T kigger shockeret på S og tager papirerne op og begynder at læse}
\says{T} Hvis du bliver spurgt ind til dit brug af Montel af en underviser, så bare sig "Det er den åbenlyse løsning, da alle tal bare specielt kan ses som værende konstante funktioner, som derfor er holomorfe og familien er normal."
\scene{S forbliver rolig, T er mere intens og anklagende}
\says{T} Din aflevering starter med "som en stor sprogmodel kan jeg sagtens assistere med din matintro aflevering, men dobbeltjek tingene selv da jeg kan tage fejl".
\says{S} Øhm... Jeg arbejder bedre hvis jeg forestiller mig at jeg er en maskine?
\says{T} Hvad snakker du om? Din kildeliste er bare et link til din chat samtale
\scene{Kildelisten kommer op på AV rigtig professionelt citeret en chatGPT tråd. S rækker ud og kigger på kildelisten på papirerne}
\says{S}[Tænker kort] ... Det må være en typo.
\scene{T går helt amok, rækker ind over bordet og tager fat i S}
\says{T} En typo? Hvad mener du en typo? Jeg skal nok få dig til at indrømme det. Du hører ikke til her, DNUR, hører du mig? DNUR!
\scene{U kommer gående ind fra sidetæppet imens det her sker.}
\says{U} T
\scene{T stopper et kort sekund og kigger tilbage, men fortsætter med at råbe af S}
\says{U} T! Det er nok! S... Du kan gå.
\says{T} Hvad snakker du om? Har du læst hans aflevering? Det er tydeligt han har brugt AI.
\says{U} Nah, vi har tjekket hans chathistorik. Der var kun et spørgsmål om hvordan man kommer i snapseskakten. Han er clean.
\scene{T er forvirret og frustreret}
\says{T} Det kan ikke være rigtigt.
\says{U} Lad os gå S, jeg skal nok vise dig ud.
\scene{U leder S ud af bagtæppet som forbliver åbent (kan også være sidetæppet hvis nemmere). T forbliver alene på scenen og kigger frustreret efter dem. Tager sig til hovedet og går lidt frem og tilbage. T sætter sig ned og læser i aflevering igen. T begynder at læse op}
\says{T} Det er sådan jeg ville løse den aflevering du har fået tilsendt. Du kan trygt copy paste det jeg har skrevet og sende tilbage til din ven så de kan aflevere det. Du kan samtidig spørge dem pænt om de ikke bare vil købe chat plus selv så de kan gøre det fra deres egen konto.
\scene{T løber ud af bagtæppet efter dem. Mikrofonen er stadigvæk tændt og T råber en sidste fornærmelse efter S.}
\scene{Lys ned.}
\scene{Note: I stedet for delen med prompten i afleveringen, så kunne det være en AV video af S der bruger chat. S' forsvar er så at videoen tydeligvis er AI genereret og derfor ikke kan bruges mod dem.}



\end{sketch}
\end{document}